\documentclass[11pt,a4paper,modern]{FFExercises}
\usepackage[english,greek]{babel}
\usepackage[utf8]{inputenc}
\usepackage{nimbusserif}
\usepackage[T1]{fontenc}
\usepackage{amsmath}
\let\myBbbk\Bbbk
\let\Bbbk\relax
\usepackage[amsbb,subscriptcorrection,zswash,mtpcal,mtphrb,mtpfrak]{mtpro2}
\usepackage{graphicx,multicol,multirow,enumitem,tabularx,mathimatika,gensymb,venndiagram,hhline,longtable,tkz-euclide,fontawesome5,eurosym,tcolorbox,tabularray,tikzpagenodes,relsize,siunitx,eurosym}
\definecolor{xrwma}{HTML}{aa1212}
\usetikzlibrary{calc}
\usetikzlibrary{positioning}
\tcbuselibrary{skins,theorems,breakable}
\renewcommand{\textstigma}{\textsigma\texttau}
\renewcommand{\textdexiakeraia}{}
\usepackage{svg}
\ekthetesdeiktes

\begin{document}

\titlos{Μαθηματικά}{Γ' Γυμνασίου}{1.9 Ρητές Αλγεβρικές Παραστάσεις}

\paragraph{ΘΕΩΡΙΑ}

\begin{enumerate}
\item \textbf{Ερωτήσεις θεωρίας}\\
Να απαντήσετε στις παρακάτω ερωτήσεις.
\begin{rlist}
\item Ποια αλγεβρική παράσταση ονομάζεται ρητή;
\item Ποια συνθήκη πρέπει να ισχύει ώστε να ορίζεται μια ρητή αλγεβρική παράσταση;
\item Πότε μπορεί να απλοποιηθεί μια ρητή αλγεβρική παράσταση;
\item Οι περιορισμοί που αφορούν μια ρητή αλγεβρική παράσταση είναι περισσότεροι ή ίδιοι μ' αυτούς που αφορούν την ίδια παράσταση απλοποιημένη;
\end{rlist}

\item \textbf{Σωστό - Λάθος}\\
Να χαρακτηρίσετε τις παρακάτω προτάσεις ως σωστές (Σ) ή λανθασμένες (Λ).
\begin{rlist}
\item Μια αλγεβρική παράσταση που περιέχει διαίρεση με μεταβλητή ονομάζεται ρητή.
\item Η παράσταση $ \dfrac{x^2+1}{3} $ είναι ρητή αλγεβρική παράσταση.
\item Οι μεταβλητές μιας ρητής παράστασης δεν μπορούν να πάρουν τιμές που μηδενίζουν τον παρονομαστή.
\item Αν σε μια ρητή παράσταση ο παρονομαστής είναι $ x^2+1 $, τότε η παράσταση ορίζεται για κάθε $ x \in \mathbb{R} $.
\item Η παράσταση $ \dfrac{2x}{x} $ είναι ίση με 2 για κάθε $ x \in \mathbb{R} $.
\item Για να απλοποιήσουμε μια ρητή παράσταση, διαιρούμε τους αριθμητές και τους παρονομαστές με τον Μ.Κ.Δ. τους.
\item Η διαίρεση δύο ρητών παραστάσεων γίνεται πολλαπλασιασμός του διαιρετέου με τον αντίστροφο του διαιρέτη.
\item Αν $ \dfrac{A}{B} = 0 $, τότε $ A=0 $ και $ B \neq 0 $.
\item Η παράσταση $ \dfrac{x-y}{y-x} $ ισούται με -1 για $ x \neq y $.
\item Το Ε.Κ.Π. των $ x^2-1 $ και $ x+1 $ είναι το $ x^2-1 $.
\item Μια ρητή αλγεβρική παράσταση ορίζεται για κάθε τιμή των μεταβλητών που περιέχει.
\item Για να απλοποιηθεί μια ρητή αλγεβρική παράσταση θα πρέπει και οι δύο όροι της να αποτελούν γινόμενο παραγόντων.
\end{rlist}
\end{enumerate}

\paragraph{ΑΣΚΗΣΕΙΣ}

\begin{center}
\textbf{ΠΕΡΙΟΡΙΣΜΟΙ}
\end{center}

\begin{enumerate}[resume]
\item \textbf{Εύρεση πεδίου ορισμού}\\
Να βρείτε τις τιμές της μεταβλητής για τις οποίες ορίζονται οι παρακάτω παραστάσεις:
\begin{multicols}{2}
\begin{rlist}
\item $ \dfrac{5}{x-3} $
\item $ \dfrac{2x}{x+4} $
\item $ \dfrac{x-1}{2x-6} $
\item $ \dfrac{3}{x^2-1} $
\item $ \dfrac{x}{x(x-2)} $
\end{rlist}
\end{multicols}

\item \textbf{Εύρεση πεδίου ορισμού (Πιο σύνθετες)}\\
Να βρείτε τις τιμές της μεταβλητής για τις οποίες ορίζονται οι παρακάτω παραστάσεις:
\begin{multicols}{2}
\begin{rlist}
\item $ \dfrac{1}{x^2-5x+6} $
\item $ \dfrac{x+1}{x^2-4x+4} $
\item $ \dfrac{2x}{x^3-4x} $
\item $ \dfrac{3}{x^2+1} $
\item $ \dfrac{x-2}{|x|-3} $
\end{rlist}
\end{multicols}

\item \textbf{Περιορισμοί}\\
Να βρεθούν οι τιμές των μεταβλητών για τις οποίες ορίζεται κάθε μια από τις παρακάτω παραστάσεις.
\begin{multicols}{3}
\begin{rlist}
\item $ \dfrac{1}{x} $
\item $ \dfrac{3}{x-1} $
\item $ \dfrac{y}{2y-4} $
\item $ \dfrac{z-3}{4-z} $
\item $ \dfrac{2x-1}{3x-6} $
\item $ \dfrac{4}{x-y} $
\end{rlist}
\end{multicols}

\item \textbf{Περιορισμοί}\\
Να βρεθούν οι τιμές των μεταβλητών για τις οποίες ορίζεται κάθε μια από τις παρακάτω παραστάσεις.
\begin{multicols}{2}
\begin{rlist}
\item $ \dfrac{3}{x^2} $
\item $ \dfrac{3x}{x(x-3)} $
\item $ \dfrac{7}{x(x+1)} $
\item $ \dfrac{x-3}{(x-2)(x+1)} $
\end{rlist}
\end{multicols}

\item \textbf{Περιορισμοί}\\
Να βρεθούν οι τιμές των μεταβλητών για τις οποίες ορίζεται κάθε μια από τις παρακάτω παραστάσεις.
\begin{multicols}{2}
\begin{rlist}
\item $ \dfrac{1}{x^2-x} $
\item $ \dfrac{2}{x^3-4x^2} $
\item $ \dfrac{3-x}{2x-x^2} $
\item $ \dfrac{x+y}{xy-y} $
\end{rlist}
\end{multicols}

\item \textbf{Περιορισμοί}\\
Να βρεθούν οι τιμές των μεταβλητών για τις οποίες ορίζεται κάθε μια από τις παρακάτω παραστάσεις.
\begin{multicols}{3}
\begin{rlist}[leftmargin=5mm]
\item $ \dfrac{x}{x^2-4} $
\item $ \dfrac{5}{4y^2-25} $
\item $ \dfrac{3-2x}{3x^2-18} $
\item $ \dfrac{3x+4}{x^3-16x} $
\item $ \dfrac{y-2}{y^2-x^2} $
\item $ \dfrac{z^2}{z^4-1} $
\end{rlist}
\end{multicols}

\item \textbf{Περιορισμοί}\\
Να βρεθούν οι τιμές των μεταβλητών για τις οποίες ορίζεται κάθε μια από τις παρακάτω παραστάσεις.
\begin{multicols}{2}
\begin{rlist}[leftmargin=5mm]
\item $ \dfrac{x+1}{x^2-4x+4} $
\item $ \dfrac{2}{x^2+2x+1} $
\item $ \dfrac{4+y}{25y^2+10y+1} $
\item $ \dfrac{z-5}{z^3-6z^2+9z} $
\end{rlist}
\end{multicols}

\begin{center}
\textbf{ΑΠΛΟΠΟΙΗΣΗ}
\end{center}

\item \textbf{Απλοποίηση Μονωνύμων}\\
Να απλοποιήσετε τις παρακάτω παραστάσεις:
\begin{multicols}{2}
\begin{rlist}
\item $ \dfrac{15x^3}{5x} $
\item $ \dfrac{12x^2y}{18xy^2} $
\item $ \dfrac{4a^2b^3}{2ab^2} $
\item $ \dfrac{5x(x-1)}{10(x-1)} $
\item $ \dfrac{3x^2(x+2)}{x(x+2)^2} $
\end{rlist}
\end{multicols}

\item \textbf{Απλοποίηση Πολυωνύμων}\\
Να παραγοντοποιήσετε αριθμητή και παρονομαστή και να απλοποιήσετε τις παραστάσεις:
\begin{multicols}{2}
\begin{rlist}
\item $ \dfrac{2x-4}{x^2-4} $
\item $ \dfrac{x^2-9}{x^2-6x+9} $
\item $ \dfrac{3x-6}{x^2-2x} $
\item $ \dfrac{x^2-1}{x^2+x} $
\item $ \dfrac{x^2-3x+2}{x^2-4} $
\end{rlist}
\end{multicols}

\item \textbf{Απλοποίηση}\\
Να απλοποιηθούν οι παρακάτω ρητές αλγεβρικές παραστάσεις.
\begin{multicols}{3}
\begin{rlist}
\item $ \dfrac{3x}{4x^2} $
\item $ \dfrac{2x}{8xy} $
\item $ \dfrac{12x^2}{4x^3y^2} $
\item $ \dfrac{5xy}{25yz} $
\item $ \dfrac{20x^2z^2}{10xz^3} $
\item $ \dfrac{8y^3z}{16yz^3} $
\end{rlist}
\end{multicols}

\item \textbf{Απλοποίηση}\\
Να απλοποιηθούν οι παρακάτω ρητές αλγεβρικές παραστάσεις.
\begin{multicols}{2}
\begin{rlist}
\item $ \dfrac{x-1}{x(x-1)^2} $
\item $ \dfrac{3x-6}{x^2-2x} $
\item $ \dfrac{4-2y}{(2-y)^2} $
\item $ \dfrac{2z+1}{4z^2-2z} $
\end{rlist}
\end{multicols}

\item \textbf{Απλοποίηση}\\
Να απλοποιηθούν οι παρακάτω ρητές αλγεβρικές παραστάσεις αφού γραφτούν οι απαραίτητοι περιορισμοί.
\begin{multicols}{2}
\begin{rlist}
\item $ \dfrac{x-3}{3x-9} $
\item $ \dfrac{2-x}{4-x^2} $
\item $ \dfrac{y^2}{y^3-y} $
\item $ \dfrac{4z^2+3z}{16z^3-9z} $
\end{rlist}
\end{multicols}

\item \textbf{Απλοποίηση}\\
Να απλοποιηθούν οι παρακάτω ρητές αλγεβρικές παραστάσεις.
\begin{multicols}{2}
\begin{rlist}
\item $ \dfrac{2x+4}{x^2+4x+4} $
\item $ \dfrac{y^2-25}{y^3-5y^2} $
\item $ \dfrac{4z^2-9}{4z^2-12z+9} $
\item $ \dfrac{t^2+10t+25}{2t^2-50} $
\end{rlist}
\end{multicols}

\item \textbf{Απλοποίηση}\\
Να απλοποιηθούν οι παρακάτω ρητές αλγεβρικές παραστάσεις.
\begin{multicols}{2}
\begin{rlist}
\item $ \dfrac{(x+y)^2}{x^2-y^2} $
\item $ \dfrac{4z-2y}{3y^2-6yz} $
\item $ \dfrac{xy^2-x^2y}{x^2-2xy+y^2} $
\item $ \dfrac{z^2-9x^2}{2zx-6x^2} $
\end{rlist}
\end{multicols}

\item \textbf{Απλοποίηση}\\
Να απλοποιηθούν οι παρακάτω ρητές αλγεβρικές παραστάσεις.
\begin{multicols}{2}
\begin{rlist}
\item $ \dfrac{x^2-3x+2}{(x-1)(x-2)} $
\item $ \dfrac{z^2-4z}{z^2-z-12} $
\item $ \dfrac{x^2-10x+25}{x^2-3x-10} $
\item $ \dfrac{y^2-5y-14}{y^2-49} $
\end{rlist}
\end{multicols}

\begin{center}
\textbf{ΣΥΝΘΕΤΕΣ ΑΣΚΗΣΕΙΣ}
\end{center}

\item \textbf{Πράξεις και Απλοποίηση}\\
Να κάνετε τις πράξεις και να απλοποιήσετε το αποτέλεσμα:
\begin{multicols}{2}
\begin{rlist}
\item $ \dfrac{x}{x+1} \cdot \dfrac{x^2-1}{x} $
\item $ \dfrac{x^2-4}{x} : \dfrac{x+2}{2x} $
\item $ \dfrac{2}{x} + \dfrac{3}{x-1} $
\item $ \dfrac{x}{x^2-1} - \dfrac{1}{x-1} $
\item $ \left( 1 + \dfrac{1}{x} \right) \cdot \dfrac{x}{x+1} $
\end{rlist}
\end{multicols}

\item \textbf{Σύνθετα Προβλήματα}\\
\begin{rlist}
\item Δίνεται η παράσταση $ A = \dfrac{x^2-16}{x^2-8x+16} $. Να βρείτε για ποιες τιμές του $ x $ ορίζεται και να την απλοποιήσετε.
\item Αν $ P(x) = \dfrac{x^3-x}{x^2+x} $, να βρείτε το πεδίο ορισμού και να δείξετε ότι $ P(x) = x-1 $.
\item Να απλοποιήσετε την παράσταση $ \dfrac{ax+ay+bx+by}{x+y} $.
\item Δίνονται τα πολυώνυμα $ A(x) = x^2-5x+6 $ και $ B(x) = x^2-9 $. Να απλοποιήσετε το κλάσμα $ \dfrac{A(x)}{B(x)} $.
\end{rlist}

\item \textbf{Πολυώνυμα - Απλοποίηση}\\
Δίνονται τα παρακάτω πολυώνυμα\\
$ \begin{aligned}
& A(x)=2x-6\\& B(x)=x^2-9\\& \varGamma(x)=x^2-7x+12
\end{aligned} $\\
Με τη βοήθεια των παραπάνω πολυωνύμων να σχηματιστούν και στη συνέχεια να απλοποιηθούν οι ακόλουθες παράστασεις :
\begin{multicols}{3}
\begin{rlist}
\item $ \dfrac{A(x)}{B(x)} $
\item $ \dfrac{B(x)}{\varGamma(x)} $
\item $ \dfrac{A(x)}{B(x)} $
\end{rlist}
\end{multicols}

\item \textbf{Ρητές - Πολυώνυμα - Τιμές - Απλοποίηση}\\
Δίνεται η ρητή παράσταση $ P(x)=\frac{A(x)}{B(x)} $ όπου οι όροι της παράστασης $ A(x), B(x) $ δίνονται από τους τύπους : \begin{gather*}
 A(x)=4-x\\ B(x)=x^2-16
\end{gather*}
\begin{rlist}
\item Να βρεθούν οι τιμές $ P(2), P(-1) $ και $ P(3) $ της παράστασης $ P(x) $.
\item Εξετάστε αν μπορεί να υπολογιστεί η τιμή $ P(4) $ της παράστασης.
\item Απλοποιήστε την παράσταση $ P(x) $ και στη συνέχεια προσπαθήστε ξανά να εξετάσετε αν υπολογίζεται η τιμή $ P(4) $. Αιτιολογίστε την απάντησή σας.
\end{rlist}

\item \textbf{Πολυώνυμα - Περιορισμοί - Απλοποίηση}\\
Με τη βοήθεια των παρακάτω πολυωνύμων\\
$ \begin{aligned}
& A(x)=4x-8\\& B(x)=x^2-4\\& \varGamma(x)=x^2-5x+6
\end{aligned} $\\
να σχηματιστούν οι ακόλουθες παράστασεις. Υπολογίστε τις τιμές της μεταβλητής για τις οποίες ορίζεται κάθε παράσταση και ύστερα απλοποιήστε τες.
\begin{multicols}{2}
\begin{rlist}
\item $ \dfrac{B(x)-A(x)}{B(x)} $
\item $ \dfrac{B(x)-\varGamma(x)}{A(x)} $
\item $ \dfrac{A(x)+\varGamma(x)}{B(x)} $
\item $ \dfrac{A(2x)}{B(-2x)} $
\end{rlist}
\end{multicols}

\end{enumerate}

\end{document}
